\documentclass[11pt]{article}

\usepackage[margin=1in]{geometry}
\usepackage{amsmath,amssymb}
\usepackage{booktabs}
\usepackage{array}
\usepackage[dvipsnames]{xcolor}
\usepackage[colorlinks=true,linkcolor=NavyBlue,urlcolor=NavyBlue]{hyperref}
\hypersetup{pdftitle={Model Guide: Bayesian Forecasting of Household Electricity Demand},
            pdfauthor={STAT 638 Project 17}}
\usepackage{enumitem}
\usepackage{fancyhdr}
\usepackage{titlesec}

\setlist{itemsep=2pt,parsep=0pt,topsep=3pt}
% Never start a page right after a displayed equation, so an equation stays
% with the explanation or symbol table that follows it, and never leave the
% first or last line of a paragraph alone on a page.
\postdisplaypenalty=10000
\clubpenalty=10000
\widowpenalty=10000
\makeatletter\@clubpenalty=\clubpenalty\makeatother
\titlespacing*{\section}{0pt}{14pt}{6pt}
\titlespacing*{\subsection}{0pt}{10pt}{4pt}

\setlength{\headheight}{14pt}
\pagestyle{fancy}
\fancyhf{}
\fancyhead[L]{\small STAT 638 Project 17 --- model guide}
\fancyhead[R]{\small\thepage}
\renewcommand{\headrulewidth}{0.4pt}

\newcommand{\kwh}{\,\mathrm{kWh}}
\newcommand{\code}[1]{\texttt{\small #1}}
\newcommand{\ci}[3]{$#1$\ \ [$#2$,\ $#3$]}

\DeclareMathOperator{\dow}{dow}
\DeclareMathOperator{\Var}{Var}
\DeclareMathOperator{\SD}{SD}
\DeclareMathOperator{\atantwo}{atan2}
\DeclareMathOperator{\GammaDist}{Gamma}
\DeclareMathOperator{\Unif}{Uniform}
\DeclareMathOperator{\Expo}{Exponential}

% Two-column reference table. It never splits across pages and stays on the
% same page as the line or equation that introduces it.
% Arguments: first-column width (fraction of \textwidth), first header, second header.
\newenvironment{reftable}[3]
  {\par\nopagebreak\addvspace{4pt}\noindent
   \renewcommand{\arraystretch}{1.3}%
   \begin{tabular}{@{}>{\raggedright\arraybackslash}p{#1\textwidth}
                      >{\raggedright\arraybackslash}p{\dimexpr0.97\textwidth-#1\textwidth-2\tabcolsep\relax}@{}}
   \toprule \textbf{#2} & \textbf{#3} \\ \midrule}
  {\bottomrule\end{tabular}\par\addvspace{8pt}}

\title{\vspace{-1.2cm}\textbf{Model Guide: Bayesian Forecasting of Household Electricity Demand}\\[2pt]
\large The model, every equation and symbol, and what the results mean}
\author{Project 17}
\date{\today}

\begin{document}
\maketitle
\thispagestyle{fancy}

\noindent
This guide collects the model-related explanations in one place: the data that
enter the model, the equations and the meaning of every symbol, the priors, how
the posterior was computed and checked, how forecasts are made and scored, what
the results say, the limitations, and the recommended next steps. Full
derivations and reproducible numbers are in \code{report.tex},
\code{src/model.py}, and \code{results/}.

% ---------------------------------------------------------------------------
\section{What is being modelled}

\begin{itemize}[leftmargin=*]
\item The data come from one household near Paris, recorded from December 2006
  through November 2010.
\item The raw file records the average active power in every minute, in
  kilowatts (kW).
\item A minute is $1/60$ of an hour, so the energy used on day $d$ is
\end{itemize}
\begin{equation}
y_d = \frac{1}{60}\sum_{m \in d} \mathrm{GAP}_m \qquad [\mathrm{kWh}].
\label{eq:agg}
\end{equation}

\begin{reftable}{0.2}{Symbol}{Meaning}
$d$ & A calendar day. \\
$m \in d$ & The recorded minutes $m$ that belong to day $d$. \\
$\mathrm{GAP}_m$ & Global active power in minute $m$, in kW. \\
$\sum_{m \in d}$ & Add up the values over all recorded minutes of day $d$. \\
$1/60$ & Converts one minute at a power of 1 kW into kWh. Equation~\eqref{eq:agg} is a daily \emph{sum}, not a daily mean. \\
$y_d$ & Total energy used on day $d$, in kilowatt-hours (kWh). \\
\end{reftable}

\begin{itemize}[leftmargin=*]
\item The two partial days at the start and end of the record are dropped. A
  remaining day counts as observed only if at least 95\% of its 1,440 minutes
  are present.
\item The result is a gap-free grid of $T = 1{,}440$ calendar days: 1,417
  observed and 23 missing. Missing days stay on the grid, so the day-to-day
  error process remains continuous.
\end{itemize}

The scientific goals are to describe long-term, yearly, and weekly patterns; to
forecast future daily demand with honest uncertainty; and to identify unusual
days or periods.

% ---------------------------------------------------------------------------
\section{The general model}

Every model fitted in the project belongs to this family:
\begin{align}
z_t &= g(y_t), \label{eq:link}\\
z_t &= \mathbf{x}_t^{\top}\boldsymbol\beta + \varepsilon_t, \label{eq:reg}\\
\mathbf{x}_t^{\top}\boldsymbol\beta
  &= \beta_0 + \beta_1\,\frac{t-\bar t}{365}
   + \sum_{j=1}^{J}\left[a_j \sin\!\left(\frac{2\pi j t}{365}\right)
   + b_j \cos\!\left(\frac{2\pi j t}{365}\right)\right] \notag\\
  &\qquad + \sum_{k=2}^{7}\delta_k\,\mathbf{1}\{\dow_t = k\}, \label{eq:mean}\\
\varepsilon_t &= \sum_{\ell=1}^{p}\phi_\ell\,\varepsilon_{t-\ell} + \eta_t,
  \qquad \eta_t \sim t_\nu(0,\sigma^2). \label{eq:ar}
\end{align}

\noindent\textbf{In words:} daily consumption $=$ a calendar-based prediction
$+$ a leftover that partly carries over from day to day, with occasional very
large surprises allowed. Equation~\eqref{eq:reg} splits each day into the
predictable part and the leftover, equation~\eqref{eq:mean} is the predictable
part, and equation~\eqref{eq:ar} describes the leftover.

\subsection{How to read the notation}

\begin{reftable}{0.25}{Notation}{Meaning}
$=$ & Equals. \\
$+$,\ $-$ & Add, subtract. \\
$a/b$ or $\frac{a}{b}$ & $a$ divided by $b$. \\
Subscript, e.g.\ $y_t$ & Identifies which day, lag, harmonic, or category the quantity belongs to. \\
Superscript 2, e.g.\ $\sigma^2$ & The quantity squared. \\
Bold, e.g.\ $\mathbf{x}$, $\boldsymbol\beta$ & A vector: a list of several numbers handled together. \\
$\sum_{j=1}^{J}$ & Add the terms for $j = 1, 2, \ldots, J$. \\
$\in$ & ``Belongs to'', e.g.\ a minute belonging to a day. \\
$\approx$ & Approximately equal to. \\
$<$,\ $>$,\ $\le$,\ $\ge$ & Less than, greater than, at most, at least. \\
$\sim$ & ``Is distributed as'': the quantity on the left follows the probability distribution on the right. \\
$\mid$ & ``Given'', or ``conditional on''. \\
$P(A \mid B)$ & Probability of $A$, given the information $B$. \\
$\Var(X)$,\ $\SD(X)$ & Variance and standard deviation of $X$. \\
$[L,\,U]$ after an estimate & Lower and upper ends of an interval. For parameters in this guide, a 95\% posterior credible interval. \\
\end{reftable}

\subsection{Response and time}

\begin{reftable}{0.2}{Symbol}{Meaning}
$t$ & Day number on the gap-free calendar, $t = 1, \ldots, T$. \\
$T$ & Number of calendar days, $T = 1{,}440$. \\
$y_t$ & Electricity used on day $t$, in kWh. Observed on 1,417 days; imputed by the model on the other 23. \\
$g(\cdot)$ & Transformation applied to $y_t$ before modelling. \\
$z_t$ & The response on the modelling scale. Identity model: $z_t = y_t$. Log model: $z_t = \log y_t$. \\
$\bar t$ & The average day index. Centring time at $\bar t$ makes the trend term zero at the middle of the record and reduces the correlation between the trend and the intercept. \\
\end{reftable}

The headline model M7 uses the identity transformation, so $z_t = y_t$ and all
of its coefficients are in kWh.

\subsection{The calendar prediction}

\begin{reftable}{0.2}{Symbol}{Meaning}
$\mathbf{x}_t$ & The known calendar inputs for day $t$: a constant 1, centred time, the sine and cosine terms, and six weekday indicators. \\
$\boldsymbol\beta$ & All regression coefficients collected into one vector. \\
$^{\top}$ & Transpose. $\mathbf{x}_t^{\top}\boldsymbol\beta$ multiplies each input by its coefficient and adds the results, giving the calendar prediction for day $t$. \\
$\beta_0$ & Intercept: expected use on a Monday at the middle of the record, before the seasonal waves are added (the waves average to zero over a year). It is not, by itself, the average-day level. \\
$\beta_1$ & Trend: change in daily use per year, because $(t-\bar t)/365$ measures time in years. \\
$365$ & Length of the yearly cycle in days, as required by the project brief. \\
\end{reftable}

The reported \emph{average-day level at the record midpoint} adds the average
weekday offset to the intercept. For M7 it is $25.96\kwh$, while the
Monday-reference intercept itself has posterior mean $23.82\kwh$.

\subsection{Annual harmonics}

\begin{reftable}{0.2}{Symbol}{Meaning}
$J$ & Number of harmonic pairs (one sine plus one cosine each). \\
$j$ & Harmonic number, $j = 1, \ldots, J$. \\
$j = 1$ & One full wave every 365 days: $\sin(2\pi t/365)$ and $\cos(2\pi t/365)$. \\
$j = 2$ & Two waves every 365 days (one every 182.5 days): $\sin(4\pi t/365)$ and $\cos(4\pi t/365)$. Added to the first harmonic, it makes the yearly curve asymmetric. \\
$a_j$ & Coefficient on the sine term of harmonic $j$, in kWh. \\
$b_j$ & Coefficient on the cosine term of harmonic $j$, in kWh. \\
$\pi$ & The constant $3.14159\ldots$ \\
$\sin(\cdot)$,\ $\cos(\cdot)$ & Smooth periodic functions. Dividing $2\pi j t$ by 365 makes harmonic $j$ repeat $j$ times a year. \\
\end{reftable}

Each sine--cosine pair can be rewritten as a single wave with a size and a
timing:
\begin{equation}
a_j\sin(\omega_j t) + b_j\cos(\omega_j t) = A_j\cos(\omega_j t - \psi_j),
\label{eq:amp}
\end{equation}
\begin{equation}
\omega_j = \frac{2\pi j}{365}, \qquad
A_j = \sqrt{a_j^2 + b_j^2}, \qquad
\psi_j = \atantwo(a_j,\, b_j).
\label{eq:phase}
\end{equation}

\begin{reftable}{0.2}{Symbol}{Meaning}
$\omega_j$ & Angular frequency of harmonic $j$, in radians per day. \\
$A_j$ & Amplitude: how far that wave rises above (and falls below) zero, in kWh. \\
$2A_j$ & Peak-to-trough swing of that wave on its own. \\
$\psi_j$ & Phase: sets the date of the wave's peak. \\
$\atantwo(a,\, b)$ & The angle whose sine and cosine are proportional to $a$ and $b$; using the signs of both puts it in the correct quadrant. \\
$\sqrt{\cdot}$ & Square root. \\
\end{reftable}

The full seasonal curve is the sum of both waves, so its total range is
generally not $2A_1 + 2A_2$.

\subsection{Weekday effects}

\begin{reftable}{0.2}{Symbol}{Meaning}
$\dow_t$ & Day of the week of day $t$ (1 = Monday, \ldots, 7 = Sunday). \\
$k$ & A weekday from Tuesday ($k = 2$) to Sunday ($k = 7$). \\
$\delta_k$ & How much weekday $k$ differs from Monday, in kWh. \\
$\mathbf{1}\{\dow_t = k\}$ & Indicator: equals 1 if day $t$ is weekday $k$, and 0 otherwise. \\
\end{reftable}

Monday has no coefficient of its own; it is the reference day. For reporting,
the seven weekday values are re-centred on their weekly mean, which answers the
more natural question: how far is this weekday above or below an average day?

\subsection{The leftover: an autoregressive error}

\begin{reftable}{0.2}{Symbol}{Meaning}
$\varepsilon_t$ & Residual (leftover) on day $t$: actual use minus the calendar prediction. \\
$p$ & Autoregressive order: how many previous days' residuals feed into today. \\
$\ell$ & Lag, $\ell = 1, \ldots, p$: how many days back. \\
$\phi_\ell$ & Persistence: the fraction of the residual from $\ell$ days ago that carries into today. \\
$\eta_t$ & Innovation: the genuinely new surprise on day $t$, after the carry-over from earlier days is accounted for. \\
\end{reftable}

For M7, $p = 1$, so the error equation is
\begin{equation}
\varepsilon_t = \phi\,\varepsilon_{t-1} + \eta_t .
\label{eq:ar1}
\end{equation}
Part of yesterday's unexplained high or low use carries into today, and then a
new surprise is added. The fitted $\phi$ is about $0.375$, so the memory fades
within a few days.

\subsection{Heavy-tailed surprises}

Each new surprise follows a Student-$t$ distribution rather than a normal one:
\begin{equation}
\eta_t \sim t_\nu(0,\, \sigma^2).
\label{eq:t}
\end{equation}

\begin{reftable}{0.2}{Symbol}{Meaning}
$t_\nu(0,\sigma^2)$ & Student-$t$ distribution centred at 0, with scale $\sigma$ and $\nu$ degrees of freedom. \\
$\sigma$ & Innovation scale: the typical size of a new one-day surprise, in kWh. \\
$\sigma^2$ & The scale squared. \\
$\nu$ & Degrees of freedom: controls how heavy the tails are. A small $\nu$ allows more extreme days; as $\nu$ grows, the distribution approaches a normal distribution. \\
\end{reftable}

Strictly, $\sigma$ is a \emph{scale}, not the standard deviation. When
$\nu > 2$,
\begin{equation}
\SD(\eta_t) = \sigma\sqrt{\frac{\nu}{\nu-2}} .
\label{eq:tsd}
\end{equation}
The project output labels $\sigma$ the ``one-step innovation sd''; ``innovation
scale'' is the precise name.

The sampler writes the Student-$t$ as a normal distribution whose variance
changes from day to day:
\begin{equation}
\eta_t \mid \lambda_t \sim N\!\left(0,\ \frac{\sigma^2}{\lambda_t}\right),
\qquad
\lambda_t \sim \GammaDist\!\left(\frac{\nu}{2},\ \frac{\nu}{2}\right).
\label{eq:mix}
\end{equation}

\begin{reftable}{0.2}{Symbol}{Meaning}
$N(\mu,\, v)$ & Normal distribution with mean $\mu$ and variance $v$. \\
$\lambda_t$ & Hidden weight for day $t$. Near 1 on an ordinary day; small on an extreme day. \\
$\sigma^2/\lambda_t$ & Variance of day $t$'s surprise. A small $\lambda_t$ makes it large, so an unusual day is absorbed without inflating the uncertainty of every ordinary day. \\
$\GammaDist(\alpha,\, \beta)$ & Gamma distribution with shape $\alpha$ and rate $\beta$. With $\alpha = \beta = \nu/2$ the weights average exactly 1. \\
\end{reftable}

Averaging over $\lambda_t$ gives back the Student-$t$ distribution in
equation~\eqref{eq:t}.

\subsection{Missing days}

The 23 missing values of $y_t$ are neither set to zero nor deleted. Their $z_t$
and $\varepsilon_t$ are sampled from the model along with everything else. This
keeps the day-to-day error link intact across gaps and carries the uncertainty
about those days into the parameter estimates and the forecasts.

% ---------------------------------------------------------------------------
\section{The headline model, M7}

The model behind the main conclusions is
\begin{align}
y_t &= \mathbf{x}_t^{\top}\boldsymbol\beta + \varepsilon_t, \label{eq:m7}\\
\mathbf{x}_t^{\top}\boldsymbol\beta
  &= \beta_0 + \beta_1\,\frac{t-\bar t}{365}
   + a_1\sin\!\left(\frac{2\pi t}{365}\right) + b_1\cos\!\left(\frac{2\pi t}{365}\right) \notag\\
  &\qquad + a_2\sin\!\left(\frac{4\pi t}{365}\right) + b_2\cos\!\left(\frac{4\pi t}{365}\right)
   + \sum_{k=2}^{7}\delta_k\,\mathbf{1}\{\dow_t = k\}, \label{eq:m7mean}\\
\varepsilon_t &= \phi\,\varepsilon_{t-1} + \eta_t,
  \qquad \eta_t \sim t_\nu(0,\sigma^2). \label{eq:m7err}
\end{align}
In short: the response is measured directly in kWh, and the model has a centred
linear trend, $J = 2$ annual harmonic pairs, six weekday indicators with Monday
as the reference, an AR(1) residual, and Student-$t$ innovations.

M7 was chosen as the headline model because it was the best of the seven
candidates on the held-out 12-month forecast, and its predictive intervals were
much sharper and better calibrated than those of the log-scale versions. It is
the best model in this candidate set, not a claim that no better model exists.

% ---------------------------------------------------------------------------
\section{Priors}

Because the model is Bayesian, every unknown parameter gets a probability
distribution before the data are seen. For M7:
\begin{align}
\beta_0 &\sim N(0,\ 100^2), \label{eq:prior-b0}\\
\beta_r &\sim N(0,\ 40^2) \quad \text{for every other regression coefficient } r, \label{eq:prior-b}\\
\sigma &\sim \text{half-}t_3(0,\ 20), \label{eq:prior-sigma}\\
\phi &\sim \Unif(-1,\ 1), \label{eq:prior-phi}\\
\nu - 2 &\sim \Expo(\text{mean } 10). \label{eq:prior-nu}
\end{align}

\begin{reftable}{0.3}{Prior}{Meaning}
$N(0,\, 100^2)$ and $N(0,\, 40^2)$ & Normal distributions centred at 0 with standard deviations 100 and 40 kWh. They are deliberately wide, so the data, not the prior, determine the coefficients. \\
$\text{half-}t_3(0,\, 20)$ & A Student-$t$ with 3 degrees of freedom and scale 20, folded at zero. It keeps $\sigma > 0$ while allowing a broad range of values. \\
$\Unif(-1,\, 1)$ & Every value of $\phi$ between $-1$ and $1$ is equally likely beforehand. $-1 < \phi < 1$ is the stationary region, where shocks fade instead of growing without bound. \\
$\Expo(\text{mean } 10)$ on $\nu - 2$ & Forces $\nu > 2$, so the variance of the Student-$t$ exists, while still allowing heavy tails. \\
Missing $y_t$ & No separate prior: their distribution is implied by the regression and the AR process. \\
\end{reftable}

The log-scale candidates use smaller prior scales because their coefficients
are much smaller: 5 for the intercept, 2 for the other coefficients, and 1 for
the half-$t$ scale.

% ---------------------------------------------------------------------------
\section{Candidate models and why they were compared}

\begin{center}
\small
\renewcommand{\arraystretch}{1.2}
\begin{tabular}{@{}lclll@{}}
\toprule
\textbf{Model} & $\boldsymbol{J}$ & \textbf{Error process} & \textbf{Scale} & \textbf{What the comparison tests} \\
\midrule
M1 & 1 & independent, Gaussian & log & naive baseline \\
M2 & 1 & AR(1), Student-$t$ & log & day-to-day memory and heavy tails \\
M3 & 2 & AR(1), Student-$t$ & log & the required second harmonic \\
M4 & 3 & AR(1), Student-$t$ & log & a more flexible yearly curve \\
M5 & 2 & AR(2), Student-$t$ & log & longer residual memory \\
M6 & 2 & AR(1), Gaussian & log & the value of heavy tails \\
\textbf{M7} & 2 & AR(1), Student-$t$ & kWh & whether the log transform helps \\
\bottomrule
\end{tabular}
\end{center}

WAIC, an in-sample estimate of predictive accuracy, preferred M5. The held-out
forecasts rejected AR(2): the extra lag was fitting long absence blocks that did
not generalise. The held-out year also showed that modelling kWh directly beat
the log transform. Forecast performance, not WAIC alone, therefore decided the
headline model.

Adding a third harmonic improved the log-scale model, but a three-harmonic
model on the kWh scale was not fitted. That remains an important extension.

% ---------------------------------------------------------------------------
\section{How the posterior was computed and checked}

The posterior distribution was computed with a purpose-built Markov chain Monte
Carlo (MCMC) sampler:
\begin{itemize}[leftmargin=*]
\item Gibbs updates for the regression coefficients, the innovation scale, the
  hidden weights $\lambda_t$, and the missing responses;
\item random-walk Metropolis updates for the AR coefficients and for
  $\log(\nu - 2)$;
\item exact stationary initialisation of the first AR residuals;
\item four chains of 40,000 iterations for the full-data fits, with the first
  10,000 of each discarded as burn-in and every fifth remaining iteration kept,
  giving 24,000 posterior draws.
\end{itemize}

\begin{reftable}{0.2}{Check}{Meaning and result}
$\widehat R$ & Compares the spread within chains with the spread between chains. Values near 1 (target below 1.01) indicate convergence. M7's worst value: 1.0069. \\
ESS & Effective sample size: how many independent draws the correlated MCMC draws are worth (target above 400). M7's smallest bulk ESS: 1,072. \\
Simulation check & On data simulated with known parameters, 14 of 15 parameters fell inside their 95\% credible intervals, as expected. \\
\end{reftable}

% ---------------------------------------------------------------------------
\section{Forecast equations}

For one posterior draw, the forecast for $h$ days after the origin $T$ is built
as
\begin{align}
\mu_{T+h} &= \mathbf{x}_{T+h}^{\top}\boldsymbol\beta, \label{eq:fc-mu}\\
\varepsilon_{T+h} &= \phi\,\varepsilon_{T+h-1} + \eta_{T+h}, \label{eq:fc-eps}\\
y_{T+h} &= \mu_{T+h} + \varepsilon_{T+h}. \label{eq:fc-y}
\end{align}

\begin{reftable}{0.2}{Symbol}{Meaning}
$T$ & Forecast origin: the last day whose data are used. \\
$h$ & Horizon: how many days ahead, $h = 1, 2, \ldots$ \\
$\mu_{T+h}$ & Calendar prediction for day $T+h$. \\
$\varepsilon_{T+h}$ & Simulated future residual, built forward one day at a time from the last known residual. \\
$y_{T+h}$ & Simulated future consumption. \\
\end{reftable}

Repeating this for every posterior draw gives the posterior predictive
distribution. It includes the uncertainty in the coefficients and in $\phi$,
$\sigma$, and $\nu$, the future innovations, and their propagation through the
AR process.

With the parameters held fixed and an innovation variance $v_\eta$, the
forecast variance of the residual $h$ days ahead is
\begin{equation}
\Var(\varepsilon_{T+h} \mid \varepsilon_T) = v_\eta \sum_{r=0}^{h-1}\phi^{2r},
\qquad
v_\eta = \sigma^2\,\frac{\nu}{\nu-2}.
\label{eq:fc-var}
\end{equation}

\begin{reftable}{0.2}{Symbol}{Meaning}
$v_\eta$ & Variance of one innovation $\eta_t$. \\
$r$ & How many days before the target day a shock occurred. \\
$\phi^{2r}$ & Contribution of that shock to the forecast variance, relative to a same-day shock. Each day of delay multiplies a shock's effect by $\phi$, and so its variance by $\phi^2$. \\
\end{reftable}

As $h$ grows, the forecast standard deviation relative to the one-day-ahead
value approaches
\begin{equation}
\frac{1}{\sqrt{1-\phi^2}} .
\label{eq:ceiling}
\end{equation}
With $\phi \approx 0.375$ the ceiling is about 1.08. Knowing the latest residual
therefore makes near-term intervals about 8\% narrower, and almost all of that
advantage is gone within four to seven days.

% ---------------------------------------------------------------------------
\section{How the forecasts were evaluated}

The final twelve months (1 December 2009 to 25 November 2010) were held out and
never used to estimate the parameters.
\begin{itemize}[leftmargin=*]
\item \textbf{Single origin:} forecast the whole held-out year from 30 November
  2009.
\item \textbf{Rolling origin:} forecast 1 to 60 days ahead from many starting
  days in the held-out year, each time updating the current residual with the
  newly observed consumption.
\end{itemize}

\begin{reftable}{0.3}{Score}{Meaning}
Log predictive density (LPD) & Rewards a forecast distribution that is both sharp and centred in the right place. Higher is better. \\
CRPS & Continuous ranked probability score: forecast error that uses the whole forecast distribution, in kWh. Lower is better. \\
Coverage & Share of actual values inside the nominal 50\% or 95\% predictive interval. Good calibration means coverage close to the nominal percentage. \\
PIT & Probability integral transform: the forecast's cumulative probability at the actual value. Roughly uniform PIT values indicate good calibration. \\
\end{reftable}

M7's rolling 95\% intervals covered 94.5\% of actual values one day ahead and
95.2\% seven days ahead. Their average width rose from about $28.9\kwh$ one day
ahead to about $31.4\kwh$ seven days ahead and then stayed almost flat.

% ---------------------------------------------------------------------------
\section{Main findings from M7}

\subsection{Overall level and trend}

All intervals in this section are 95\% posterior credible intervals.

\begin{reftable}{0.5}{Quantity}{Posterior mean [95\% interval]}
Average-day level at the record midpoint & \ci{25.96}{25.41}{26.50} kWh/day \\
Trend & \ci{-0.186}{-0.667}{0.292} kWh/year \\
Trend as a percentage of the level & \ci{-0.72}{-2.57}{1.13} \% per year \\
$P(\text{slope} < 0 \mid \text{data})$ & $0.775$ \\
\end{reftable}

The interval for the trend includes zero, so there is no conclusive long-term
rise or fall.

\subsection{Yearly pattern}

\begin{reftable}{0.5}{Quantity}{Posterior mean [95\% interval]}
First-harmonic amplitude $A_1$ & \ci{8.08}{7.31}{8.85} kWh \\
First-harmonic peak-to-trough swing $2A_1$ & \ci{16.15}{14.62}{17.71} kWh \\
First-harmonic peak date & about 14 January\ \ [9 January, 20 January] \\
Second-harmonic amplitude $A_2$ & \ci{2.43}{1.68}{3.19} kWh \\
Amplitude ratio $A_2/A_1$ & \ci{0.302}{0.204}{0.405} \\
$P(A_2 > 0.15\,A_1 \mid \text{data})$ & $0.999$ \\
Full seasonal range (maximum minus minimum) & \ci{18.25}{16.48}{20.06} kWh \\
\end{reftable}

The second harmonic clearly matters: it sharpens the winter peak and makes the
summer low longer and flatter.

\subsection{Weekly pattern}

Deviations from the weekly mean, in kWh per day:

\begin{reftable}{0.5}{Day}{Posterior mean [95\% interval]}
Monday & \ci{-2.15}{-2.91}{-1.40} \\
Tuesday & \ci{-0.41}{-1.17}{0.37} \\
Wednesday & \ci{-0.10}{-0.89}{0.69} \\
Thursday & \ci{-2.75}{-3.51}{-1.99} \\
Friday & \ci{-1.12}{-1.88}{-0.36} \\
Saturday & \ci{+3.50}{2.67}{4.33} \\
Sunday & \ci{+3.02}{2.21}{3.85} \\
Weekend minus weekday & \ci{4.56}{3.71}{5.42} \\
$P(\text{weekend} > \text{weekday} \mid \text{data})$ & $1.0000$ \\
\end{reftable}

\subsection{Error process}

\begin{reftable}{0.5}{Quantity}{Posterior mean [95\% interval]}
Innovation scale $\sigma$ & \ci{5.59}{5.25}{5.94} kWh \\
AR coefficient $\phi$ & \ci{0.375}{0.323}{0.426} \\
Degrees of freedom $\nu$ & \ci{5.81}{4.38}{7.81} \\
Long-run to one-step forecast sd ratio, $1/\sqrt{1-\phi^2}$ & \ci{1.079}{1.057}{1.105} \\
\end{reftable}

These values show modest short-term persistence and clearly heavier tails than
a normal distribution.

\subsection{Unusual days and periods}

The day-by-day check uses the posterior predictive p-value
\begin{equation}
p_t = P\!\left(y_t^{\mathrm{rep}} \le y_t \mid \mathbf{y}\right).
\label{eq:ppp}
\end{equation}

\begin{reftable}{0.2}{Symbol}{Meaning}
$y_t^{\mathrm{rep}}$ & A replicate of day $t$ simulated from the fitted model. \\
$\mathbf{y}$ & The full observed data set. \\
$p_t$ & Probability that a replicate is no larger than what actually happened. Near 0: unusually low use. Near 1: unusually high use. \\
\end{reftable}

\begin{itemize}[leftmargin=*]
\item The one-day-ahead check flagged 31 days in the extreme 1\% tails: 23
  unusually high and 8 unusually low, against about 14 of each expected by
  chance.
\item A one-day-ahead check can miss a sustained absence, because the AR(1)
  term adapts after the first day or two.
\item A separate check on 14-day average consumption found 13 sustained unusual
  episodes.
\item The clearest low periods look like summer vacancies; the high winter
  episodes are consistent with cold spells.
\end{itemize}
These are interpretations based on timing and load shape, not on measured
occupancy or weather.

% ---------------------------------------------------------------------------
\section{What the model tells us in practice}

\begin{itemize}[leftmargin=*]
\item Winter versus summer is the largest predictable difference.
\item This household uses more electricity on weekends than on weekdays.
\item There is not enough evidence to plan around a long-term rise or fall.
\item Yesterday's consumption improves forecasts mainly for the next few days.
\item Beyond about a week, the forecast is effectively a seasonal-plus-weekday
  distribution.
\item Single-day forecasts stay wide because one household's behaviour is
  inherently variable.
\item Plan with forecast intervals; a point forecast alone hides most of the
  relevant uncertainty.
\end{itemize}

% ---------------------------------------------------------------------------
\section{Limitations}

\begin{itemize}[leftmargin=*]
\item This is one household, not a population or an electricity grid. The
  results should not be generalised without new data.
\item There is no temperature variable. The harmonics capture average weather
  for the time of year but cannot predict a specific cold snap.
\item There is no occupancy variable. Vacations are treated as outliers rather
  than as a separate away-from-home state.
\item Weekday effects are additive and constant, although the exploratory work
  shows a larger weekend premium in winter.
\item The innovation scale is constant, although residual variation is larger
  in winter than in summer.
\item Trend and season are mildly confounded, because the record is just short
  of four complete years.
\item The required 365-day period ignores leap-year drift.
\item Conditional WAIC is not fully reliable for choosing forecasting models
  with autocorrelated data; the held-out results are more trustworthy.
\item A simple day-of-year average (``climatology'') beat every Bayesian
  candidate on the single-origin one-year forecast. Part of its advantage comes
  from similar August absences in the training and test years, but it also
  shows that a two-harmonic yearly curve is not flexible enough for the best
  long-range forecasts.
\end{itemize}

% ---------------------------------------------------------------------------
\section{Recommended next steps}

\begin{enumerate}[leftmargin=*]
\item Add daily Paris-area temperature; this is the most valuable missing
  predictor.
\item Fit kWh-scale models with more harmonics, $J = 3$ through $J = 6$, since
  only the log-scale three-harmonic model was tested.
\item Compare versions with and without the linear trend for long-horizon
  forecasts.
\item Let the innovation scale vary by season, so winter can be more variable
  than summer.
\item Let the weekend effect depend on the season.
\item Consider a hidden home/away state for multi-week absences.
\item Include the day-of-year climatology in the same rolling-origin evaluation
  as the Bayesian candidates.
\item Use a rolling origin on every day, rather than every third day, and score
  all horizons on a common set of target dates.
\end{enumerate}

% ---------------------------------------------------------------------------
\section{Distinctions that are easy to mix up}

\begin{reftable}{0.3}{Distinction}{Explanation}
Coefficient vs.\ reported effect & $\delta_k$ is measured relative to Monday; the reported weekday profile is re-centred on the weekly mean. \\
Residual vs.\ innovation & $\varepsilon_t$ includes carry-over from earlier days; $\eta_t$ is only the new one-day shock. \\
Credible vs.\ predictive interval & A credible interval describes uncertainty about a parameter; a predictive interval describes uncertainty about a future observation and is much wider. \\
Innovation scale vs.\ overall variability & $\sigma$ governs one new shock; AR persistence makes longer-run residual variability larger. \\
In-sample fit vs.\ forecast skill & WAIC summarises fit to the data used for estimation; the held-out year measures forecasting directly. \\
Best candidate vs.\ best possible model & M7 is the headline model among M1--M7; the climatology result and the untested kWh models show clear room for improvement. \\
\end{reftable}

% ---------------------------------------------------------------------------
\section{Where to verify each part}

\begin{reftable}{0.45}{Topic}{File}
Core equations, sampler, forecasting, diagnostics & \code{src/model.py} \\
Candidate models and WAIC & \code{src/03\_fit\_models.py} \\
Held-out evaluation & \code{src/04\_holdout.py} \\
Derived effects and anomaly checks & \code{src/05\_derived.py} \\
Posterior summaries & \code{results/03\_convergence.txt} \\
WAIC comparison & \code{results/03\_waic.txt} \\
Forecast results & \code{results/04\_holdout.txt} \\
Interpreted posterior quantities & \code{results/05\_derived.txt} \\
Full technical narrative & \code{report.tex} \\
\end{reftable}

\end{document}
